\chapter{Introduzione}\label{ch:introduzione}

\section{Descrizione}
\textbf{DriveHub} è un'applicazione desktop per le attività essenziali di un concessionario: presentare i veicoli disponibili, organizzare test drive e noleggi, gestire le richieste di acquisto e valutare le auto proposte dai clienti. Un unico catalogo collega queste operazioni allo stato di ciascun veicolo, così che le scelte degli utenti si riflettano sulla sua disponibilità.

Il sistema prevede tre figure operative:
\begin{itemize}
  \item \textbf{Cliente (\code{CLIENT}):} si registra, consulta il catalogo, prenota un test drive o un noleggio, versa l'acconto per richiedere un acquisto e propone al concessionario un proprio veicolo. Può consultare le proprie operazioni e usare un conto dimostrativo per i movimenti previsti dall'applicazione.
  \item \textbf{Venditore (\code{SALESMAN}):} esamina i veicoli proposti dai clienti e decide se accettarli o rifiutarli; autorizza il completamento delle richieste di acquisto e può cambiare la destinazione di un veicolo tra vendita e noleggio. Nella propria area consulta anche richieste, noleggi e test drive.
  \item \textbf{Manager (\code{MANAGER}):} cura il catalogo inserendo marche, modelli e veicoli, applica sconti a modelli o singole auto e consulta le vendite registrate.
\end{itemize}

Le tre aree condividono le informazioni sul veicolo, ma espongono operazioni diverse in base al ruolo. Nel codice il cliente è identificato da \code{CLIENT}; la denominazione \emph{Customer} compare anche nei materiali di progetto.

\section{Stack tecnologico}
\begin{itemize}
  \item \textbf{Linguaggio e interfaccia:} Java 21 e JavaFX 21.0.2. I layout sono definiti in FXML, lo stile in CSS e le azioni delle schermate nei controller.
  \item \textbf{Organizzazione del codice:} i package \code{presentation}, \code{business}, \code{domain} e \code{dao} separano interfaccia, casi d'uso, regole del dominio e accesso ai dati. Le interfacce DAO sono distinte dalle implementazioni JDBC.
  \item \textbf{Persistenza:} PostgreSQL tramite driver JDBC. All'avvio, se la connessione non riesce, l'applicazione usa un database H2 in memoria in modalità compatibile con PostgreSQL, inizializzato con dati dimostrativi.
  \item \textbf{Sicurezza:} le password sono salvate come hash derivati con PBKDF2, anziché in chiaro.
  \item \textbf{Build e verifica:} Maven gestisce dipendenze e compilazione; JUnit 5 e Mockito supportano i test, JaCoCo la misurazione della copertura.
\end{itemize}

\section{Semplificazioni adottate}\label{sec:semplificazioni}
Per mantenere il progetto realizzabile e dimostrabile, alcuni processi del concessionario sono rappresentati in forma ridotta:
\begin{itemize}
  \item \textbf{Pagamenti interni:} acconti, saldi, noleggi e rimborsi modificano il saldo di un \code{BankAccount} locale. Non sono integrati banche o circuiti di pagamento e i movimenti non costituiscono transazioni finanziarie reali.
  \item \textbf{Account dello staff:} la registrazione consente di scegliere anche i ruoli Venditore e Manager. È una scelta del prototipo; l'attribuzione di questi ruoli in un sistema operativo richiederebbe un controllo amministrativo.
  \item \textbf{Modalità dimostrativa:} H2 permette di avviare l'applicazione senza un server PostgreSQL, ma i dati in memoria non sopravvivono alla chiusura. Questa modalità serve a provare i flussi principali.
  \item \textbf{Processi aziendali:} il sistema registra le operazioni sul veicolo senza modellare sedi, fornitori, officina, consegna fisica, contratti, fatturazione o assistenza successiva alla vendita.
\end{itemize}

La documentazione in \code{docs/} conserva anche requisiti e risultati di una baseline precedente: fra questi, ordini di stock, un gateway di pagamento con stati propri e prove PostgreSQL/GUI più estese. I capitoli 2--4 descrivono il checkout presente in \code{src/} e in \code{src/main/resources/db/migration/init.sql}; quando richiamano un requisito storico non realizzato, ne dichiarano esplicitamente il limite. I risultati dei test sono attribuiti soltanto alla baseline su cui sono stati eseguiti.

Il progetto ha impiegato un assistente di intelligenza artificiale per analisi comparativa, revisione di testo, diagrammi, codice e test; il processo e le verifiche dichiarate dagli autori sono registrati in \code{docs/USO_AI.md}. La presente revisione ha confrontato i file locali e rieseguito la suite standard. I capitoli successivi presentano requisiti, progettazione, implementazione e verifiche del sistema.
